% TOP_PROBLEM: {"id": 30002589, "title": "Deligne's Conjecture for Critical Motive Values", "category_id": 1, "category": {"id": 1, "name": "number_theory", "display_name": "Number Theory", "description": "Properties of integers, prime numbers, Diophantine equations.", "slug": "number-theory", "order_index": 1, "created_at": "2026-07-31T15:26:25.670Z"}, "status": "partially_solved"}
% FIRST_POSTED: 2026-09-27
% !TeX program = lualatex
% Compile twice: lualatex 30002589.tex
% Standalone research notebook; original section preserved.
% Research additions: 27 September 2026. No external TeX input or BibTeX file.
\documentclass[11pt,a4paper]{article}
\usepackage[margin=24mm,headheight=14pt]{geometry}
\usepackage{fontspec}
\setmainfont{Latin Modern Roman}
\setsansfont{Latin Modern Sans}
\setmonofont{Latin Modern Mono}
\usepackage{amsmath,amssymb,mathtools}
\usepackage{unicode-math}
\setmathfont{Latin Modern Math}
\usepackage{microtype}
\usepackage{enumitem,xurl,needspace}
\usepackage[dvipsnames]{xcolor}
\usepackage{hyperref}
\hypersetup{unicode=true,colorlinks=true,linkcolor=MidnightBlue,urlcolor=MidnightBlue,
 pdftitle={Research notebook - Problem 30002589},pdfauthor={Open Problems Math},bookmarksnumbered=true}
\usepackage{bookmark}
\usepackage{titlesec}
\titleformat{\section}{\Large\bfseries\raggedright}{\thesection}{0.7em}{}
\usepackage{fancyhdr}
\pagestyle{fancy}\fancyhf{}
\fancyhead[L]{\small\sffamily Open Problems Math | Research notebook}
\fancyhead[R]{\small\sffamily Problem 30002589 | 27 September 2026}
\fancyfoot[C]{\thepage}
\setlength{\parindent}{0pt}
\setlength{\parskip}{5pt plus 1pt}
\setlength{\emergencystretch}{3em}
\setcounter{tocdepth}{1}
\setcounter{secnumdepth}{1}
\setlist[itemize]{leftmargin=3em,itemsep=5pt,topsep=4pt,parsep=0pt}
\allowdisplaybreaks
\numberwithin{equation}{section}
\newcommand{\N}{\mathbb{N}}
\newcommand{\Z}{\mathbb{Z}}
\newcommand{\Q}{\mathbb{Q}}
\newcommand{\R}{\mathbb{R}}
\newcommand{\C}{\mathbb{C}}
\newcommand{\F}{\mathbb{F}}
\newcommand{\A}{\mathbb{A}}
\newcommand{\PP}{\mathbb P}
\newcommand{\E}{\mathbb{E}}
\newcommand{\eps}{\varepsilon}
\newcommand{\dd}{\,\mathrm d}
\newcommand{\sref}[2]{\hyperlink{source:#1:#2}{[S#2]}}
\newcommand{\eref}[2]{\hyperlink{extra:#1:#2}{[E#2]}}
\newif\ifshowcatalogue
\showcataloguetrue
\newcommand{\cataloguescope}[1]{\ifshowcatalogue
 \paragraph{Exact scope in the supplied catalogue.}
 \begin{quote}\small #1\end{quote}\fi}
\begin{document}
\setcounter{section}{0}
\section*{\texorpdfstring{Deligne's Conjecture for Critical Motive Values}{Deligne's Conjecture for Critical Motive Values}}\label{30002589}
\subsection{Definitions and mathematical statement}
For a pure motive $M/\Q$ with coefficient field $E$, let $H_B^\pm(M)$ be the complex-conjugation eigenspaces of its Betti realization and $d_\pm$ their $E$-dimensions. The periods $c^\pm(M)$ are determinants of the specified comparison maps to de Rham filtration quotients, defined modulo $E^\times$. For a critical integer $n$, neither archimedean factor for $M$ at $n$ nor for $M^\vee$ at $1-n$ has a pole. Put $s=(-1)^n$. The assertion is
\[
 \frac{L(M,n)}{(2\pi i)^{n d_s}c^s(M)}\in E
 \ \subseteq\ E\otimes_{\Q}\C\simeq\prod_{\sigma:E\hookrightarrow\C}\C.
\]
The membership is diagonal across all embeddings. The finite $L$-values and compatible realizations are prerequisites; continuation and nonvanishing are not additional conclusions of this period law. The recorded scope fixes the period and sign conventions.
\par\smallskip\textit{Formulation and concept references:} \sref{30002589}{1}.
\cataloguescope{For every pure motive M over Q with coefficients in a number field E, in Deligne's geometric realization formalism (Betti, de Rham and compatible l-adic realizations with comparison, dual, tensor and Tate-twist structures), and every critical integer n, the normalized finite L-value vector L(M,n)/((2*pi*i)\textasciicircum{}(n*d\_sign)*c\textasciicircum{}sign(M)) lies in the diagonal E inside E tensor C, where sign=(-1)\textasciicircum{}n and d\_sign=dim\_E H\_B\textasciicircum{}sign(M). Here H\_B\textasciicircum{}sign is the sign eigenspace of complex conjugation, and c\textasciicircum{}sign is the determinant of the Betti/de Rham comparison to the corresponding filtration quotient, computed in E-rational bases and defined modulo E-units as in Deligne1.7/2.6. Critical means that L\_infinity(sigma,M,s) and L\_infinity(sigma,Mdual,1-s) have no pole at s=n for every embedding sigma:E->C (equivalently use Mdual(1) at -s). Use finite, not completed, L-functions with the compatible E-valued local factors of2.2. Analytic continuation giving finite values at the selected critical integer and the source's well-defined compatible realizations/local factors are defining prerequisites, not assertions proved by this period law and not a restriction to cases whose continuation is already known. The simultaneous embedding vector must come from one element of E; unrelated componentwise algebraicity is insufficient. All-zero values are allowed, but mixed zero/nonzero components do not satisfy the diagonal membership. No general nonvanishing or order-of-zero assertion from Deligne2.7 is appended.}
\subsection{Short English statement}
After normalization by the predicted periods, do all coefficient-field embeddings of a critical motivic $L$-value arise from one element of that field?
% BEGIN RESEARCH_ADDITION ATTEMPT_71
\subsection{Research attempt}

\paragraph{Current frontier and normalization audit.}
Deligne's original Sections~1.7 and~2.6 were inspected, including the
comparison determinants and the coefficient-field formulation
\eref{30002589}{1}. The supplied sign and finite-$L$ conventions are retained.
Recent factorization results for automorphic motives over CM fields have
regularity and archimedean rationality hypotheses; they do not establish
the statement for every pure motive over $\Q$ \eref{30002589}{2}.
The present attempt makes no new analytic-continuation or nonvanishing
assertion: the selected finite values and compatible realizations remain
prerequisites as in the original record.

\paragraph{Attack route.}
Restriction of coefficients can turn an $E$-valued assertion into scalar
statements, but may keep only a norm and lose the compatibility between
embeddings. Componentwise algebraicity has the same weakness. The proposed
repair is to use all coefficient-weighted traces, or a finite set of
coefficient-weighted determinant probes. The linear-algebraic reduction
will be proved exactly; the missing motivic step is to realize those probes
by comparison identities, rather than treating arbitrary scalar operations
on $L$-values as new motives.

\paragraph{Attempt: trace-dual recovery of a diagonal element.}
Let $d=[E:\Q]$ and write $B=E\otimes_{\Q}\C$. In its product description,
$z=(z_\sigma)_{\sigma:E\hookrightarrow\C}$ is an arbitrary vector; in the
application it is the period-normalized vector in the original statement.
Define the $\C$-linear functional
$\operatorname{Tr}_B(z)=\sum_\sigma z_\sigma$.
Choose a $\Q$-basis $e_1,\ldots,e_d$ of $E$ and its trace-dual basis
$e_1^*,\ldots,e_d^*$, characterized by
$\operatorname{Tr}_{E/\Q}(e_i e_j^*)=\delta_{ij}$.
This basis exists because the trace pairing of a number field is
nondegenerate. Extending that pairing to $\C$ gives
\begin{equation}\label{eq71:reconstruct}
 z=\sum_{i=1}^d\operatorname{Tr}_B(e_i z)e_i^*.
\end{equation}
Indeed, both sides have the same pairing with every $e_j$, which is a
$\C$-basis of $B$. Nondegeneracy then gives equality. In particular,
\begin{equation}\label{eq71:criterion}
 z\in E\quad\Longleftrightarrow\quad
 \operatorname{Tr}_B(e_i z)\in\Q\quad(1\leq i\leq d).
\end{equation}
No prior algebraicity of the individual $z_\sigma$ is needed. The converse
direction is not a heuristic use of all embeddings: the displayed formula
constructs the one element of $E$ they must represent.

Applied to
\[
 z_\sigma=\frac{L(M,\sigma,n)}
 {(2\pi i)^{n d_s}c^s(M)_\sigma},\qquad s=(-1)^n,
\]
this gives a finite family of scalar rationality targets. It preserves the
source's sign and comparison determinants; it is not a replacement of them
by their absolute values. A change of allowed $E$-rational bases multiplies
$z$ by a diagonal element of $E^\times$. Since multiplication by such an
element maps $E$ bijectively onto itself, the criterion is independent of
that choice, exactly as the original period law requires \eref{30002589}{1}.

\paragraph{Attempt: recovering the traces from finite determinant probes.}
For each basis element form the degree-at-most-$d$ polynomial
\[
 Q_i(T)=\operatorname{Norm}_{B/\C}(1+T e_i z)
       =\prod_\sigma(1+T\sigma(e_i)z_\sigma).
\]
Its linear coefficient is $\operatorname{Tr}_B(e_i z)$.
Fix $d+1$ distinct rational numbers $t_0,\ldots,t_d$. If
\begin{equation}\label{eq71:probes}
 Q_i(t_j)\in\Q\qquad(1\leq i\leq d,
                                     \ 0\leq j\leq d),
\end{equation}
then all coefficients of $Q_i$ are rational. To see this, the coefficient
vector is obtained from its value vector by inverting the rational
Vandermonde matrix $(t_j^k)_{0\leq j,k\leq d}$. Its determinant is
$\prod_{a<b}(t_b-t_a)\ne0$, so its inverse also has rational entries.
Thus \eqref{eq71:probes} implies \eqref{eq71:criterion} and hence $z\in E$.
Conversely, $z\in E$ makes every $Q_i$ a rational polynomial and implies
all these probe identities. This is a finite exact equivalence, not an
approximation by numerically recognized algebraic numbers.

The useful conceptual difference from a lone norm test is that multiplication
by the $e_i$ remembers which coefficient embedding is which. A hypothetical
coefficient-compatible comparison construction for these probes would
settle the period membership assertion without requiring separate
nonvanishing. At present no such general construction is supplied. In
particular, tensor products multiply local Euler factors in prescribed
ways; they do not make $1+t e_i z$ into the normalized critical value of
an automatically available auxiliary motive.

\paragraph{Stress test: a counterexample to scalar-only descent.}
Take $E=\Q(\sqrt2)$ and order the two embeddings by
$\sqrt2\mapsto\sqrt2,-\sqrt2$. Set $z=(2,3)$.
Both coordinates are rational, its trace is $5$, and its norm is $6$.
Even the whole characteristic polynomial
\[
 (T-2)(T-3)=T^2-5T+6
\]
is rational. Nevertheless $z\notin E$: if
$(a+b\sqrt2,a-b\sqrt2)=(2,3)$ with $a,b\in\Q$, then
$a=5/2$ and $b=-1/(2\sqrt2)\notin\Q$.
The omitted probe detects the defect:
\[
 \operatorname{Tr}_B(\sqrt2 z)=-\sqrt2\notin\Q.
\]
Thus all characteristic coefficients of the unweighted vector being
rational is still weaker than the target. Numerical matches of a norm,
individual minimal polynomials, or independent algebraicity in each
embedding do not supply the common coefficient-field element.

Likewise $(0,1)$ cannot be a diagonal element of any quadratic field:
an embedding of a field is injective, so an element with one zero embedding
is zero in every embedding. In contrast, the zero vector satisfies
\eqref{eq71:criterion} and \eqref{eq71:probes}; the argument introduces no
extraneous nonvanishing demand. These examples concern arbitrary vectors
in $B$, not actual motivic $L$-values, and are therefore not counterexamples
to Deligne's conjecture. The exact verification script checks the quadratic
basis and probes symbolically.

\paragraph{Outcome.}
\textbf{Outcome: Conditional reduction.}

The coefficient-equivariant conclusion is reduced exactly to finitely many
weighted trace rationalities, equivalently to a finite rational determinant
probe test. The reconstruction and the failure of unweighted scalar
information are proved. These are elementary descent lemmas developed for
the notebook, without a priority claim.

The decisive missing input is arithmetic: establish those weighted
identities for the actual normalized comparison vector of every motive
and every critical integer in the given scope. Scalar special-value results
for restricted automorphic families do not supply that input automatically.
No new general period theorem, analytic continuation, or nonvanishing
statement is claimed.
% END RESEARCH_ADDITION ATTEMPT_71
\subsection{Sources}
\begin{itemize}\raggedright
\item[\textnormal{[S1]}]\hypertarget{source:30002589:1}{} Kazuki Morimoto, On special values of L-functions for quaternion unitary groups of degree2 and GL(2), Oberwolfach Report22/2014; with Deligne's definitions and twist convention.
\url{https://ems.press/content/serial-article-files/46514?nt=1}.
{\small\textit{Catalogue formulation source.}}
\item[\textnormal{[S2]}]\hypertarget{source:30002589:2}{} Factorization of periods, construction of automorphic motives and Deligne's conjecture over CM-fields.
\url{https://arxiv.org/abs/2509.02303}.
{\small\textit{Catalogue research/status reference; not a proof endorsement.}}
\item[\textnormal{[S3]}]\hypertarget{source:30002589:3}{} Motivic pieces of curves: L-functions and periods.
\url{https://arxiv.org/abs/2601.21934}.
{\small\textit{Catalogue research/status reference; not a proof endorsement.}}
% BEGIN RESEARCH_ADDITION REFERENCES_71
\item[\textnormal{[E1]}]\hypertarget{extra:30002589:1}{} Pierre Deligne,
\emph{Valeurs de fonctions $L$ et p\'eriodes d'int\'egrales},
Proceedings of Symposia in Pure Mathematics 33, part 2 (1979), 313--346,
Sections 1.7 and 2.1--2.6.
\url{https://publications.ias.edu/sites/default/files/33_Valeursde.pdf}.
{\small\textit{Original comparison and coefficient conventions inspected
27 September 2026; supplements the incorporated definitions in [S1].}}
\item[\textnormal{[E2]}]\hypertarget{extra:30002589:2}{} Harald Grobner,
Michael Harris and Jie Lin, \emph{Factorization of periods, construction of
automorphic motives and Deligne's conjecture over CM-fields},
arXiv:2509.02303 (2025), abstract and stated automorphic hypotheses.
\url{https://arxiv.org/abs/2509.02303}.
{\small\textit{Inspected 27 September 2026; regularity and archimedean
rationality assumptions are not removed.}}
% END RESEARCH_ADDITION REFERENCES_71
\end{itemize}

\end{document}
